\documentclass[12pt]{article}
\usepackage[margin=1in]{geometry}
\usepackage{amsmath,amssymb,amsthm}
\usepackage{algorithm}
\usepackage{algpseudocode}
\usepackage{bm}
\newtheorem{proposition}{Proposition}

\title{A No-Wait Single-Intersection Scheduler by Priority Enumeration:\\
A Decision Tree with a Longest-Path Timing Model}
\author{}\date{}

\begin{document}
\maketitle

\section{Problem Formulation}\label{sec:form}

Let $\mathcal{N}=\{1,\dots,N\}$ index the vehicles to be scheduled through the
merging zone, and let $\mathcal{M}$ be the set of conflict points (\emph{spaces}).
Vehicle $n$ traverses an ordered sequence of sub-tasks
$s\in\mathcal{S}_n=\{1,\dots,S_n\}$; sub-task $(n,s)$ occupies conflict point
$\sigma(n,s)\in\mathcal{M}$ for a fixed duration $c_{n,s}>0$. Define the prefix
offset and the total traversal length
\begin{equation}
\pi_{n,s}=\sum_{k=1}^{s-1}c_{n,k},\qquad
C_{n}=\sum_{s\in\mathcal{S}_n}c_{n,s},
\end{equation}
and let $\alpha_n$ denote the earliest admissible entry (release) time of
vehicle $n$.

\subsection{No-wait motion}
Each vehicle crosses the intersection without stopping, so its sub-tasks are
contiguous. The schedule of vehicle $n$ is therefore parameterized by a
\emph{single} entry time $t_n\ge\alpha_n$, under which sub-task $(n,s)$ occupies
$\sigma(n,s)$ over the interval
\begin{equation}
I_{n,s}(t_n)=\bigl[\,t_n+\pi_{n,s},\;\; t_n+\pi_{n,s}+c_{n,s}\,\bigr].
\end{equation}

\subsection{Scheduling problem}
Two sub-tasks conflict iff they use the same conflict point; collect such pairs in
\begin{equation}
\mathcal{K}=\bigl\{\,\{(n,s),(n',s')\}:\sigma(n,s)=\sigma(n',s'),\;n\neq n'\,\bigr\}.
\end{equation}
A schedule is \emph{feasible} iff conflicting occupations never overlap.
Minimizing the total delay yields
\begin{equation}\tag{P}\label{eq:P}
\begin{aligned}
\min_{t\in\mathbb{R}^{N}}\;\; & \sum_{n\in\mathcal{N}}\bigl(t_n-\alpha_n\bigr)\\
\text{s.t.}\;\; & t_n\ge\alpha_n,\quad n\in\mathcal{N},\\
& I_{n,s}(t_n)\cap I_{n',s'}(t_{n'})=\varnothing,
\;\;\forall\{(n,s),(n',s')\}\in\mathcal{K}.
\end{aligned}
\end{equation}
Since motion is no-wait, the exit time of $n$ is $t_n+C_n$, so
$\sum_n(t_n-\alpha_n)$ equals the total completion delay.

\section{Disjunctive Form and the Selection}\label{sec:disj}

Each non-overlap constraint is a disjunction. For
$\{(n,s),(n',s')\}\in\mathcal{K}$, the orientation \emph{``$(n,s)$ before
$(n',s')$''} reads
\begin{equation}\label{eq:arc}
t_{n'}-t_{n}\;\ge\;w_{n\to n'}\;:=\;\pi_{n,s}+c_{n,s}-\pi_{n',s'},
\end{equation}
and the reverse orientation has weight
$w_{n'\to n}=\pi_{n',s'}+c_{n',s'}-\pi_{n,s}$.
A \emph{selection} $\xi$ orients every conflict pair, i.e.\ it is a set of
weighted precedence arcs $i\xrightarrow{\,w\,}j$ on the vehicle nodes.
Adjoining the release arcs $0\xrightarrow{\,\alpha_n\,}n$ from a source node $0$
yields the precedence graph $G_\xi=\bigl(\{0\}\cup\mathcal{N},\,A_\xi\bigr)$.

\section{Timing Model}\label{sec:timing}

\begin{proposition}[Earliest-entry timing]\label{prop:timing}
If $G_\xi$ is acyclic, then problem \eqref{eq:P} restricted to the orientation
$\xi$ has a unique pointwise-minimal solution $t^\star(\xi)$, characterized by
the longest-path recursion
\begin{equation}\label{eq:timing}
\boxed{\;
t^\star_n(\xi)=\max\!\Bigl(\alpha_n,\;
\max_{\,i:\,(i\to n)\in A_\xi}\bigl(t^\star_i(\xi)+w_{i\to n}\bigr)\Bigr),
\quad n\in\mathcal{N},\;}
\end{equation}
computed by the fixed-point \textnormal{(}Bellman--Ford\textnormal{)} iteration
$\mathsf{EarliestT}(\xi)$. If $G_\xi$ contains a positive cycle, $\xi$ is
infeasible.
\end{proposition}

Equation \eqref{eq:timing} is the \emph{timing model}
$\mathsf{EarliestT}:\xi\mapsto t^\star(\xi)$: it maps an ordering to the unique
contention-free, no-wait schedule that starts every vehicle as early as the
chosen precedences permit. Its cost is
\begin{equation}\label{eq:J}
J(\xi)\;=\;\sum_{n\in\mathcal{N}}\bigl(t^\star_n(\xi)-\alpha_n\bigr).
\end{equation}

\section{Reduction to a Finite Search}\label{sec:reduction}

\begin{proposition}\label{prop:reduce}
$\displaystyle \min\eqref{eq:P}=\min_{\xi\in\Xi}J(\xi)$, where $\Xi$ is the
finite set of acyclic selections.
\end{proposition}

The optimum is attained by a semi-active schedule; hence it suffices to search
over orderings, the timing model \eqref{eq:timing} supplying the optimal entry
times of each.

\section{Decision-Tree Search}\label{sec:tree}

A node of the search tree is a \emph{partial} selection $\hat\xi$ with decided
set $D\subseteq\mathcal{K}$. The search interleaves the timing model
\eqref{eq:timing} with conflict detection. The conflict oracle
$\mathsf{FindConflict}(t,D)$ returns the first pair
$\{(n,s),(n',s')\}\in\mathcal{K}\setminus D$ whose intervals overlap at the
current schedule $t$, i.e.\ satisfying
\begin{equation}\label{eq:overlap}
t_n+\pi_{n,s}<t_{n'}+\pi_{n',s'}+c_{n',s'}
\;\;\wedge\;\;
t_{n'}+\pi_{n',s'}<t_n+\pi_{n,s}+c_{n,s},
\end{equation}
and $\textsc{none}$ if no such pair exists.

\begin{algorithm}[t]
\caption{No-wait enumeration scheduler}\label{alg:tree}
\begin{algorithmic}[1]
\Procedure{Recurse}{$\hat\xi,\,D$}
  \State $t\gets\mathsf{EarliestT}(\hat\xi)$
         \Comment{timing model, Prop.~\ref{prop:timing}}
  \If{$t=\textsc{none}$} \State \Return
         \Comment{positive cycle $\Rightarrow$ infeasible (prune)}
  \EndIf
  \State $\kappa\gets\mathsf{FindConflict}(t,D)$
         \Comment{first overlapping, undecided pair \eqref{eq:overlap}}
  \If{$\kappa=\textsc{none}$}
     \State record leaf $t^\star\gets t$; \Return
         \Comment{conflict-free complete schedule}
  \EndIf
  \State $\{(n,s),(n',s')\}\gets\kappa$
  \State \Call{Recurse}{$\hat\xi\cup\{\,n\!\xrightarrow{w_{n\to n'}}\!n'\,\},\;
         D\cup\{\kappa\}$}
         \Comment{$(n,s)$ before $(n',s')$}
  \State \Call{Recurse}{$\hat\xi\cup\{\,n'\!\xrightarrow{w_{n'\to n}}\!n\,\},\;
         D\cup\{\kappa\}$}
         \Comment{$(n',s')$ before $(n,s)$}
\EndProcedure
\State \textbf{Initialize:} \Call{Recurse}{$\varnothing,\,\varnothing$}
\end{algorithmic}
\end{algorithm}

The set of leaves enumerates $\{t^\star(\xi):\xi\in\Xi\}$, i.e.\ all
contention-free complete schedules; the minimum-cost leaf solves \eqref{eq:P}.
A positive-cycle test in $\mathsf{EarliestT}$ prunes infeasible orientations, so
only acyclic selections are expanded.

\end{document}
